\chapter{Introduction}
\section{Introduction}With more than a billion speakers and growing pull in global business, tech, and culture, Mandarin isn't just a language, it's basically a key to what's coming. China's economic rise, its grip on manufacturing and AI, plus its widening cultural reach mean that learning Mandarin can open real career doors. But for a lot of learners, the complexity, those tones, those tangled characters, makes actually mastering it feel hopeless. Here's what the research on language learning tells us: people pick things up best through meaningful exposure to authentic content~\citep{krashen1985,Mayer_2009}. By authentic, I mean real stuff made for native speakers, books, movies, news, podcasts, conversations, not the dumbed-down materials built for learners. When you're engaging with things you honestly care about, motivation climbs and you remember more. It mirrors how kids learn their first language, honestly. Through immersion, not memorization~\citep{cychosz2021, gowenlock2024}. Traditional methods? They trip a lot of students up. Textbooks and classrooms lean hard on grammar rules rather than natural exposure to the language. Apps these days hand you translations but stumble over Mandarin's quirks. The character system throws people because the written symbols barely hint at how they sound. And Pinyin helps with pronunciation, sure, though it rarely shows up anywhere outside a classroom. LexiFlow\footnote{The live application is accessible at \url{https://lexiflow.amerai.top}, and the source code is available on GitHub at \url{https://github.com/amer-adam/lexiflow}.} tackles all this with a web platform that pairs sentence segmentation with machine translation models. It works through audio by spotting sentence boundaries and treating each chunk on its own. Doing it that way bumps up accuracy and cuts down on the computing power needed. The system also keeps a database of stuff it's already processed, so there's no duplicate work, and learners get materials that match their interests and where they're at.
\section{Problem Background}Language lecturers at UTM have noticed something off in how their students develop. Plenty of them ace the written exams by memorizing vocabulary and grammar rules, but the moment they're dropped into a real Mandarin conversation, they freeze up. That gap, between scoring well on paper and actually being able to talk, gets in the way of real learning. Authentic content fixes a lot of this, since it exposes students to useful words, the way people actually speak, and cultural context no textbook really captures, and it works even better when it lines up with what students personally care about. But making that content? Still a slog. Lecturers start by hunting down real audio or video, which alone eats up a surprising amount of time. Once they've picked something, they sit down and transcribe the spoken Mandarin by hand, add pinyin, then write out the translation. For just a few minutes of material this can run 8 to 10 hours, and that's a heavy load to put on teachers. The tools out there don't really help. Take LanguageReactor, it leans completely on YouTube subtitles existing, and those are pretty rare for Mandarin compared to English. Then there's stuff like Transcribe.com, which spits out content without the timestamped subtitle format students need to match a sound to its written form, plus no pinyin, which Mandarin learners kind of depend on for pronunciation. What students actually need is a system that chews through authentic Mandarin content and shows the spoken language right next to its written version. Accurate timestamped transcriptions, pinyin guides, solid translations, all tying natural speech back to the text.
\section{Project Aim}Build LexiFlow, a web platform that handles, transcribes, and translates real Mandarin using smart sentence-level segmentation and translation models. Basically to help students actually get how real-world Mandarin conversations work.
\section{Project Objectives}\label{sec:obj}
Here's what the project's trying to do:
\begin{enumerate}[label= (\alph*)] \item Figure out what the principal stakeholders actually need, language learners and educators included. \item Build a web platform with a user interface and a database that serves up synchronized, time-stamped bilingual subtitles. \item Develop that platform around what stakeholders asked for, throwing in features like audio/video processing, transcription, translation, and pinyin generation. Plus \item Test and validate it through proper testing methods and feedback from domain experts, so the transcription, translation, and subtitle timing all stay accurate.
\section{Project Scopes}The scopes of the project are:
\begin{enumerate}[label= (\alph*)] \item The platform takes long-form audio or video, up to 120 minutes, either from files users upload or straight from YouTube URLs. \item It runs voice activity detection (VAD) along with end-of-sentence detection, breaking the audio into sentence-level chunks that actually make linguistic sense. \item Transcriptions come from automatic speech recognition (ASR), while translations are handled by a Neural Machine Transformer (NMT) model. \item Using the transformer model, the platform builds synchronized, timestamped bilingual subtitles, original text, the target-language translation, and pinyin annotations all together. \item Processed subtitles get stored in a searchable database indexed by content identifiers (think URLs), so the same content can be reused instantly without reprocessing. \item And there's a web interface where users upload content and tweak how the subtitles come out.
\section{Project Importance}LexiFlow matters because it's really about making language learning better. The big problem? Online platforms and lecturers only supply you so much. LexiFlow fixes that. By pulling in machine translation models for both translation and transcription, it lets students figure out from content they actually care about.
\section{Report Organization}This thesis breaks down into six chapters. Chapter 1 kicks things off, laying out the background, the purpose, the objectives, the scope, why it matters, and how the report's put together. Chapter 2 digs into the literature, looking at existing systems and tech around speech recognition, translation models, and the stuff you have to think about with user interfaces, plus a comparison against what's already on the market. Chapter 3 walks through how the system was actually built. It covers the software development life cycle I went with, why I picked that technique and the main stages along the strategy. Chapter 4 is all about requirements and design, digging into functional and non-functional requirements, the database setup, and the interface itself, so the whole thing stays usable and in fact works. Chapter 5, Implementation and Testing, gets into how it was actually built and verified: the architecture in practice, the core features, deployment, and both automated and user acceptance testing. The final one, Chapter 6, wraps everything up. It pulls together the results, talks through whether the project hit its goals, and suggests where LexiFlow could go next. The notion behind this layout is simple: a clear path from concept all the way to implementation and evaluation.
